\medskip

	En France il y a deux formules pour obtenir le permis de conduire :

	\begin{itemize}[label*={\textbullet~}]
		\item Suivre à partir de 15 ans une formation de conduite accompagnée pendant 2 ans ;
		\item Suivre la formation classique (sans conduite accompagnée) à partir de 17 ans.
	\end{itemize}

	\medskip

	En France actuellement, parmi les jeunes qui suivent une formation au permis de conduire, 16~\% choisissent la formation de conduite accompagnée, et parmi eux, 74,7~\% réussissent l'examen de conduite dès leur première tentative.


	En suivant la formation classique, le taux de réussite dès la première tentative est seulement de 56,8~\%.

	\bigskip

	On choisit au hasard un jeune français qui a déjà passé l'examen de conduite et on considère les évènements $A$ et $R$ suivants :

	\begin{itemize}[label*={\textbullet~}]
		\item $A$ : \og le jeune a suivi la formation de conduite accompagnée\fg{};
		\item $R$ : \og le jeune a eu le permis dès sa première tentative \fg{}.
	\end{itemize}

	\medskip

	\textbf{On arrondira les résultats à $\mathbf{10}^{\mathbf{-3}}$ près, si nécessaire.}

	\medskip

	\subsection*{Partie A}
		\begin{enumerate}
			\item Dresser un arbre de probabilités modélisant cette situation.
			\item \begin{enumerate}
				\item Démontrer que $P(R)= \np{0,59664}$.
			\end{enumerate}

			Dans la suite, on gardera la valeur 0,597 arrondie à $10^{-3}$ près.

			\begin{enumerate}[resume]
				\item Donner ce résultat en pourcentage et l'interpréter dans le contexte de l'exercice.
			\end{enumerate}

			\item On choisit un jeune ayant eu son permis dès sa première tentative. Quelle est la probabilité qu'il ait suivi la formation de conduite accompagnée ?

			\item Quelle devrait être la proportion de jeunes suivant la formation de conduite accompagnée si on voulait que le taux de réussite global (quelle que soit la formation choisie) dès la première tentative à l'examen de conduite dépasse $70~\%$ ?
		\end{enumerate}

	\subsection*{Partie B}
		Une auto-école présente pour la première fois à l'examen de conduite 10 candidats qui ont suivi la formation de conduite accompagnée. On modélise le fait de passer les examens de conduite par des épreuves aléatoires indépendantes.

		\medskip

		On note $X$ la variable aléatoire donnant le nombre de ces 10 candidats qui auront leur permis dès la première tentative.


		\begin{enumerate}
			\item Justifier que $X$ suit une loi binomiale de paramètres $n=10$ et $p=0,747$.

			\item Calculer $P(X \geqslant 6)$. Interpréter ce résultat.

			\item Déterminer $E(X)$ et $V(X)$.

			\item Il y a aussi 40 candidats qui n'ont pas suivi la formation de conduite accompagnée et qui se présentent pour la première fois à l'examen de conduite. De la même manière, on note $Y$ la variable aléatoire qui donne le nombre de ces candidats qui auront le permis à la première tentative. On admet que $Y$ est indépendante de la variable $X$ et qu'en fait $E(Y)=22,53$ et $V(Y)=9,81$.

			\medskip

			On note alors $Z$ la variable aléatoire comptant le nombre total de candidats (parmi les 50) qui auront le permis de conduire dès la première tentative dans cette auto-école.

			\medskip

			\begin{enumerate}
				\item Exprimer $Z$ en fonction de $X$ et $Y$. En déduire $E(Z)$ et $V(Z)$.

				\item En utilisant l'inégalité de Bienaymé-Tchebychev, montrer que la probabilité qu'il y ait au plus 20 ou au moins  40 candidats qui aient leur permis dès la première tentative est inférieure à 0,12.
			\end{enumerate}
		\end{enumerate}

\bigskip
